Game design innovation is constrained by designers' awareness of the existing design
space. This paper applies k-means clustering ($k = 6$) to 20 games embedded in the
five-dimensional TIGER fingerprint space and identifies four sparse regions ---
combinations of TIGER scores that are computationally distant from all corpus games.
The analysis reveals Codenames as a singleton cluster with no near neighbours at
distance $< 2.5$, confirming architectural isolation, and surfaces a surprising
structural equivalence between Pandemic and Chess (shared cluster, centroid
$(T, I, G, E, R) = (7.1, 6.8, 4.6, 7.3, 6.7)$). Four design gap briefs derived from
sparse cluster coordinates are validated by a panel of five experienced game designers:
three of four gaps ($75\%$) are rated as commercially viable ($\geq 3.5/5$) and
genuinely unexplored by $\geq 3$ of 5 designers, surpassing the pre-registered threshold
of $\geq 2$ validated gaps. The one non-validated gap (deep-forgiving Euro) is rejected
because designers identify existing games in that region not captured by the corpus,
revealing a Euro-tradition sampling bias. The paper demonstrates TIGER's utility as a
generative design tool: not merely for evaluating games that exist, but for discovering
where new games could be made.
